\emph{(i).} On vectors $M(F)a = a\cdot F = \tfrac12(aF - Fa)$ is the map with matrix $F^{\mu}{}_{\nu}$, antisymmetric with respect to the metric, so $\mathrm{tr}\,M(F)^{2} = F^{\mu}{}_{\nu}F^{\nu}{}_{\mu} = -F_{\mu\nu}F^{\mu\nu}$, and with the normalisation $\langle F^{2}\rangle_{0} = -X/2$ of \claimref{synthesis.one-loop-running-coefficients-from-spacetime-algebra-spin-traces} this is $-X$: $s = -4$ in every signature and dimension. On a left ideal $AP$, $P$ an idempotent, the real trace of left multiplication by $x$ is the trace over the whole algebra of $y \mapsto xyP$, which is $\sum_{b}\langle b^{-1}xbP\rangle_{0}$ over the blades $b$. For a blade $x$ the sum $\sum_{b}b^{-1}xb$ is $2^{d}x$ if $x$ is central and $0$ otherwise (each generator either commutes or anticommutes with $x$, and a non-central blade anticommutes with half of the blades). $F^{2} = \langle F^{2}\rangle_{0} + \langle F^{2}\rangle_{4}$, and blades of grade four are not central for $d \le 5$; so the real trace of $\tfrac14F^{2}$ on $AP$ is $\tfrac14\langle F^{2}\rangle_{0}$ times the real dimension $2n_{c}$, its complex trace $\tfrac14 n_{c}\langle F^{2}\rangle_{0} = -n_{c}X/8$, and $s = -n_{c}/2$. The check evaluates both traces on a generic rational bivector in each algebra, with ideals of $n_{c} = 2$ complex components in dimension three and $4$ in four and five.

\emph{(ii).} By clause (iii) of Fact \claimref{synthesis.imported-heat-kernel-coefficients-and-background-field-gauge} the quantum gauge field is a real boson on the vector index ($\sigma = 1$, $w = \tfrac12$, $n = d$, $s = -4$): $\tfrac13\cdot\tfrac12(d - 24)$; its ghosts are a complex scalar of opposite statistics ($\sigma = -1$, $w = 1$, $n = 1$, $s = 0$): $-\tfrac13$; the squared Dirac operator has $\sigma = -1$, $w = \tfrac12$, $n = n_{c}$, $s = -n_{c}/2$: $-\tfrac16(n_{c} - 3n_{c}) = n_{c}/3$. By clause (i) of the Fact the flat local part of $a_{4}$ in dimension $n$ is $(4\pi)^{-n/2}\int\mathrm{tr}(\tfrac12E^{2} + \tfrac1{12}\Omega^{2})$ and enters $\mathrm{Tr}\,e^{-tD}$ with $t^{(4-n)/2}$; the proper-time integral $\int dt\,t^{-1}t^{(4-n)/2}$ is logarithmic at its lower end only for $n = 4$, converges there for $n = 3$ and diverges as a power for $n = 5$. The gauge field has mass dimension $(d - 2)/2$ with the coupling in the covariant derivative, so the squared coupling has dimension $4 - d$.

\emph{(iii).} Three dimensions: by \claimref{dark-tower-matter.three-dimensional-content}(i) each factor couples to one two-component fermion of unit charge; an abelian factor has no adjoint Casimir. Five dimensions: by \claimref{dark-tower-matter.five-dimensional-content} the species under (S) are $U$ $(3, \tfrac23)$, $D$ $(3, -\tfrac13)$, $E$ $(1, -1)$, $N$ $(1, 0)$, four-component Dirac fields, and without (S) one triplet of $Y = \tfrac23$; $C_{2} = 3$ and $T = \tfrac12$ on the triplet (clause (3) of the spin-trace result); $\sum\dim\cdot Y^{2} = \tfrac43 + \tfrac13 + 1 = \tfrac83$. The four-dimensional reading replaces $-\tfrac72$ by $-\tfrac{11}{3}$: $-11 + \tfrac43 = -\tfrac{29}{3}$. The algebras are those of \claimref{synthesis.gauge-algebra-by-tower}.

\emph{(iv).} Stored statements: \claimref{dark-matter.mirror-tower-copy-one}(iii), (iv) for the mirror and for the scope of the common-coupling premise; \claimref{dark-tower-matter.three-dimensional-phases} for the phases of the three-dimensional towers. The one-loop relation $\Lambda = \mu_{0}\exp(2\pi/(b\,\alpha_{0}))$ for $b < 0$ gives the displayed difference with $b = -7$ and $-\tfrac{29}{3}$. The clause is a reading: it lists what a ratio would need and finds none of it stored.
